% 4_system_design.tex

\section{Visual Analytics Design}
\label{sec:vis-design}

The dashboard follows a coordinated multiple view (CMV) architecture~\cite{roberts2007coordinated} in which spatial, analytical, and temporal perspectives on EVCP siting are tightly linked through shared state. The design is guided by the visual analytics mantra~\cite{keim2008visual}---\emph{analyse first, show the important, zoom, filter, analyse further, details on demand}---and draws on established principles for multi-criteria spatial decision support~\cite{andrienko2007geovisual,jankowski2001gis}.

%--------------------------------------------------------------
\subsection{Dashboard Layout}
\label{sec:layout}

The interface is organised into three coordinated panels (Figure~\ref{fig:dashboard-layout}):

\begin{enumerate}[nosep]
    \item \textbf{Weight Laboratory} (left panel): interactive controls for MCDA method selection, criterion weighting via parallel coordinates, pairwise AHP comparison, and a weight diagnostics view.
    \item \textbf{Spatial View} (centre panel): a web map displaying H3 hexagonal choropleth and glyph overlays, with a DFES timeseries panel anchored below.
    \item \textbf{Impact Panel} (right panel): scenario placement list, eight KPI cards, a radar chart for scenario comparison, and a per-LSOA breakdown of aggregated impacts.
\end{enumerate}

\noindent All panels share state through a reactive store architecture (Zustand), ensuring that weight changes, map interactions, and scenario edits propagate immediately across views without explicit event wiring.

\begin{figure}[htbp]
    \centering
    % PLACEHOLDER: Screenshot of the full dashboard showing all three panels
    \fbox{\parbox{0.95\columnwidth}{\centering\vspace{3cm}\textbf{[Figure: Dashboard Layout]}\\\small Full dashboard view showing the Weight Laboratory (left), Spatial View with H3 choropleth (centre), and Impact Panel with KPI cards (right).\vspace{0.5cm}}}
    \caption{The dashboard layout with three coordinated panels. The Weight Laboratory provides interactive weight controls and AHP pairwise comparison. The Spatial View renders MCDA suitability scores as a hexagonal choropleth with optional glyph overlays. The Impact Panel displays KPIs and scenario comparison charts.}
    \label{fig:dashboard-layout}
\end{figure}

%--------------------------------------------------------------
\subsection{Weight Laboratory}
\label{sec:weight-lab}

The Weight Laboratory provides two complementary mechanisms for criterion weight assignment:

\paragraph{Direct manipulation via parallel coordinates.}
A parallel coordinate display renders the $m$ active criteria as vertical axes. Each axis carries a draggable handle that controls the criterion weight; dragging upward increases the weight while the remaining weights are proportionally adjusted to maintain the unit-sum constraint. A horizontal bar chart below the parallel coordinates provides a secondary encoding of the weight distribution. Criteria can be toggled on or off, and their benefit/cost polarity can be switched, with each change triggering an immediate MCDA re-computation.

\paragraph{Pairwise comparison via AHP.}
An interactive AHP matrix interface (Figure~\ref{fig:ahp-interface}) enables structured weight elicitation. Users adjust pairwise importance ratios on the Saaty scale ($\tfrac{1}{9}$ to $9$) using sliders. The system solves the comparison matrix in real time using the Row Geometric Mean Method (Section~\ref{sec:ahp}), displaying the derived weights, the Consistency Ratio (CR), and a consistency status indicator. When $\text{CR} < 0.10$, the derived weights can be applied to the MCDA model with a single action, replacing the current direct-manipulation weights. A diagnostics panel shows the full comparison matrix, enabling users to identify and correct inconsistent judgements.

\begin{figure}[htbp]
    \centering
    % PLACEHOLDER: Screenshot of the AHP comparison interface
    \fbox{\parbox{0.95\columnwidth}{\centering\vspace{3cm}\textbf{[Figure: AHP Interface]}\\\small The AHP pairwise comparison panel showing sliders for each criterion pair, derived weights, and the Consistency Ratio indicator.\vspace{0.5cm}}}
    \caption{The AHP pairwise comparison interface. Sliders control the relative importance of each criterion pair. The derived weight vector and Consistency Ratio are computed and displayed in real time.}
    \label{fig:ahp-interface}
\end{figure}

\paragraph{Method selection.}
A dropdown control allows switching between WSM, WPM, and TOPSIS aggregation methods. Changing the method triggers a new DuckDB query with the same weight vector but a different scoring function, enabling rapid comparison of how methodological choice affects the spatial ranking pattern.

%--------------------------------------------------------------
\subsection{Spatial View and Glyph Encoding}
\label{sec:spatial-view}

The primary spatial encoding is a hexagonal choropleth rendered as GeoJSON polygons over a CartoDB Positron basemap using MapLibre GL JS~\cite{maplibregl}. Each H3 cell is coloured by its MCDA suitability score using a perceptually uniform \texttt{viridis} colour scale, with opacity adjustable via a continuous slider.

An optional glyph overlay (Figure~\ref{fig:glyph-map}) projects per-cell criterion profiles as either \emph{bar glyphs} (up to 6 criteria) or \emph{rose glyphs} (up to 8 criteria), drawn on HTML5 Canvas and positioned at cell centroids. The glyph layer encodes multiple visual variables following Bertin's~\cite{bertin1983semiology} framework:

\begin{itemize}[nosep]
    \item \textbf{Size}: petal length (rose) or bar height encodes the normalised criterion value.
    \item \textbf{Width}: petal angular sweep (rose only) encodes the criterion weight, making higher-weighted criteria visually dominant.
    \item \textbf{Texture}: a diagonal hatch pattern distinguishes \emph{cost} criteria from \emph{benefit} criteria within the same glyph, following the convention that texture signals qualitative difference.
    \item \textbf{Comparison mode}: clicking a glyph pins it as a reference; subsequent hover glyphs show a dashed outline of the pinned profile for direct pairwise comparison.
\end{itemize}

\noindent The glyph aggregation resolution is user-controlled (H3~r6--r10), enabling analysts to switch between local detail and regional overview without recomputing MCDA scores.

\begin{figure}[htbp]
    \centering
    % PLACEHOLDER: Screenshot showing glyph overlay on map
    \fbox{\parbox{0.95\columnwidth}{\centering\vspace{3cm}\textbf{[Figure: Glyph Map]}\\\small Map view with rose glyph overlay showing per-cell criterion profiles. Cost criteria are distinguished by a hatch pattern. A pinned reference glyph (solid) enables comparison with a hovered glyph (dashed outline).\vspace{0.5cm}}}
    \caption{The glyph overlay showing rose glyphs at H3 resolution~8. Each petal represents a criterion: length encodes the normalised value, angular width encodes the weight, and hatching distinguishes cost criteria. The comparison mode overlays a pinned reference profile (solid) against the hovered cell (dashed).}
    \label{fig:glyph-map}
\end{figure}

%--------------------------------------------------------------
\subsection{Simulation Mode and Placement Interaction}
\label{sec:simulation-mode}

Activating \emph{Simulation Mode} changes the map cursor to an EVCP placement marker and enables a click-to-place interaction on any H3 cell. On click, the system:

\begin{enumerate}[nosep]
    \item Retrieves the MCDA query result row for the clicked cell, extracting both raw and normalised criterion values as well as LSOA and borough metadata from the underlying DuckDB tables.
    \item Opens a placement configuration popover displaying the reverse-geocoded address (via Nominatim), a charger type selector, a charger count slider, and the estimated installation cost.
    \item On confirmation, stores the placement with its full spatial context---criterion values, LSOA code, LSOA name, and borough name---ensuring that impact calculations remain reproducible even if the user later changes MCDA weights.
\end{enumerate}

\noindent This \emph{capture-at-placement-time} design is a deliberate architectural choice: by snapshotting the spatial data when the placement is created, the impact model is decoupled from the current MCDA state, preventing silent invalidation of impact estimates when criterion weights change.

%--------------------------------------------------------------
\subsection{Impact Panel: KPIs and Radar Comparison}
\label{sec:impact-panel}

The right panel (Figure~\ref{fig:impact-panel}) renders eight KPI cards (Section~\ref{sec:demand}) with large-format numbers and contextual labels. When multiple scenarios are saved, a radar chart normalises six axes---energy, carbon, revenue, population, utilisation, and headroom spare---to $[0, 1]$ using the maximum value across all series, enabling direct visual comparison of scenario profiles.

\begin{figure}[htbp]
    \centering
    % PLACEHOLDER: Screenshot of Impact Panel with KPI cards and radar chart
    \fbox{\parbox{0.95\columnwidth}{\centering\vspace{3cm}\textbf{[Figure: Impact Panel]}\\\small The Impact Panel showing KPI cards (top), a radar chart comparing two saved scenarios (middle), and a per-LSOA breakdown table (bottom).\vspace{0.5cm}}}
    \caption{The Impact Panel. KPI cards display energy delivered, carbon saved, installation cost, revenue, peak demand, headroom impact, population served, and utilisation. The radar chart enables multi-dimensional comparison of saved scenarios.}
    \label{fig:impact-panel}
\end{figure}

\subsubsection{Per-LSOA Breakdown}

For multi-placement scenarios, a collapsible section groups placements by their LSOA code (carried as metadata on each placement) and reports per-group summaries including charger count, aggregated energy, carbon, peak demand, and borough name. This enables planners to assess spatial concentration and identify areas receiving disproportionate or insufficient investment.

%--------------------------------------------------------------
\subsection{Scenario Contextualisation via DFES Projections}
\label{sec:dfes}

A key analytical contribution is the integration of scenario-level impact estimates with UK Power Networks' Distribution Future Energy Scenarios (DFES)~\cite{ukpn-dfes}, transforming the DFES timeseries from a passive reference panel into a coordinated supply-versus-demand visualisation.

\subsubsection{Data and Geographic Scope}

DFES projections are retrieved at runtime from the UKPN Open Data portal via a paginated REST API. Each record associates a year, a scenario pathway (\emph{Leading the Way}, \emph{Consumer Transformation}, \emph{System Transformation}, \emph{Falling Short}), a geographic identifier, and projected adoption counts for four technology metrics: electric cars, electric vans, heat pumps, and domestic batteries. The panel supports three levels of geographic scoping---Network, Borough, and LSOA---selectable via a toggle control. When the user clicks an H3 cell on the map, the LSOA code and borough name propagate to the DFES panel, which auto-selects the appropriate scope.

\subsubsection{Supply--Demand Overlay}

The scenario's aggregated impact KPIs are converted into DFES-commensurate units. For example, annual energy delivered is converted to the number of EVs supportable ($V_{\text{EV}} = E_{\text{annual}} / 2{,}500~\text{kWh}$), and peak demand to equivalent domestic battery installations ($V_{\text{BAT}} = P_{\text{peak}} / 4~\text{kW}$). The converted value is rendered as a horizontal dashed line across the DFES timeseries chart.

Shaded fill areas between the supply line and each demand pathway curve communicate the supply--demand relationship (Figure~\ref{fig:dfes-overlay}):

\begin{itemize}[nosep]
    \item \textbf{Green fill} (below the supply line): projected demand is within the scenario's supply capacity.
    \item \textbf{Red fill} (above the supply line): projected demand exceeds the supply capacity, indicating unmet demand.
\end{itemize}

\subsubsection{Year of Constriction}

For each DFES pathway, the platform identifies the \emph{year of constriction}---the first year in which the projected demand curve crosses the supply line from below, computed via linear interpolation:

\begin{equation}
\label{eq:constriction}
t^{*} = t_i + \frac{V - y_i}{y_{i+1} - y_i} \cdot (t_{i+1} - t_i)
\end{equation}

\noindent This is encoded as a filled circle at the intersection point with a vertical drop-line to the $x$-axis, providing a direct temporal answer to the planning question: ``Under which pathway, and in which year, will this deployment become insufficient?''

\begin{figure}[htbp]
    \centering
    % PLACEHOLDER: Screenshot of DFES timeseries with supply-demand overlay
    \fbox{\parbox{0.95\columnwidth}{\centering\vspace{3cm}\textbf{[Figure: DFES Overlay]}\\\small DFES timeseries showing four demand pathways (coloured lines) with the scenario supply line (dashed horizontal). Green shading indicates sufficient capacity; red shading indicates unmet demand. Year-of-constriction markers appear at pathway--supply intersections.\vspace{0.5cm}}}
    \caption{The DFES timeseries panel with supply--demand overlay. The horizontal dashed line represents the scenario's supply capacity (converted to DFES metric units). Shaded regions encode the temporal supply--demand relationship. Intersection markers indicate the year of constriction for each demand pathway.}
    \label{fig:dfes-overlay}
\end{figure}

%--------------------------------------------------------------
\subsection{Scenario Management}
\label{sec:scenario-mgmt}

The Scenario Manager enables users to save the complete analytical state---MCDA weights, method, criterion polarities, AHP comparisons, and all placements---as a named scenario. Saved scenarios can be restored, compared side-by-side via the radar chart, imported/exported as JSON files, and toggled between visible, highlighted, and hidden display modes. This supports iterative planning workflows where a planner evaluates multiple deployment strategies and selects the preferred alternative based on multi-dimensional performance comparison.

%--------------------------------------------------------------
\subsection{Coordinated Linking Model}
\label{sec:linking}

Interactions propagate across views through a shared reactive store:

\begin{itemize}[nosep]
    \item \textbf{Map $\rightarrow$ DFES}: clicking an H3 cell sets the selected LSOA and borough; the DFES panel re-fetches data scoped to the selected geography.
    \item \textbf{Map $\rightarrow$ Impact}: placing a charger triggers impact recalculation; KPI cards and the radar chart update within the same render cycle.
    \item \textbf{Weight Lab $\rightarrow$ Map}: adjusting a criterion weight re-runs the DuckDB MCDA query; the choropleth and glyph overlays update after a 200\,ms debounce to maintain responsive interaction.
    \item \textbf{DFES $\rightarrow$ Map}: selecting an LSOA in the DFES panel updates the map's selected geography and adjusts the scope accordingly.
\end{itemize}

\noindent This architecture supports \emph{loose coupling}: each panel subscribes only to the state slices it needs, minimising unnecessary re-renders during high-frequency interactions such as weight dragging.
